\documentclass{article}
\usepackage{graphicx} % Required for inserting images
\usepackage{amssymb}

\title{Chapter 29 Problem Set}
\author{Siddharth K}
\date{May 2024}

\begin{document}

\maketitle

\textbf{1a.} The sequence $a_n$ converges to $\frac{3}{5}$. Given $\epsilon > 0$, let $N = \frac{1}{25\epsilon} - \frac{7}{5}$. Then $|\frac{3n + 4}{5n + 7} - \frac{3}{5}| = |\frac{15n - 20 - 15n - 21}{25n + 35}| = \frac{1}{25n + 35} < \frac{1}{25(\frac{1}{25\epsilon} - \frac{7}{5}) + 35} = \epsilon$.
\newline

\textbf{1b.} The sequence converges to $0$. We see that $|\frac{n}{n^2 + 3}| < \epsilon$ when $\epsilon n^2 - n + 3\epsilon > 0$. Since $\epsilon > 0$, this is an upward opening parabola. This means, at some point (which we can call $N$), all the values of $a_n$ are greater than $0$. So $|\frac{n}{n^2} + 3| < \epsilon$ when $n > N$.
\newline

\textbf{1c.} The sequence diverges. Given $\epsilon > 0$, we need to prove that there exists $N$ such that if $n > N$ then $\frac{n^2 - 2}{n + 1} > \epsilon$. That is true if $n^2 - \epsilon n - \epsilon - 2 > 0$. Since that is an upward opening parabola, that is true for all $n$ greater than some point, $N$.
\newline

\textbf{1d.} The sequence converges to $0$. If $N = \frac{1}{\epsilon}$, then $|\frac{(-1)^n}{n}| = |\frac{1}{n}| < \frac{1}{N} = \epsilon$.
\newline

\textbf{1e.} The sequence converges to $0$. If $N = \frac{1}{\epsilon}$, then $|\frac{\sin n}{n}| \le |\frac{1}{n}| < \frac{1}{N} = \epsilon$.
\newline

\textbf{1f.} The sequence converges to $0$. If $N = \frac{1}{\epsilon}$, then $|\frac{\cos n}{n}| \le |\frac{1}{n}| < \frac{1}{N} = \epsilon$.
\newline

\textbf{1g.} The sequence converges to $\frac{1}{2}$. If $N = \frac{\epsilon}{2} + \frac{1}{2} + \frac{1}{8\epsilon}$, then $$\left|n - \sqrt{n^2 - n} - \frac{1}{2}\right| < \left|\frac{\epsilon}{2} + \frac{1}{2} + \frac{1}{8\epsilon} - \sqrt{\left(\frac{\epsilon}{2} + \frac{1}{2} + \frac{1}{8\epsilon}\right)^2 - \left(\frac{\epsilon}{2} + \frac{1}{2} + \frac{1}{8\epsilon}\right)} - \frac{1}{2}\right|$$$$= \left|\frac{\epsilon}{2} + \frac{1}{2} + \frac{1}{8\epsilon} - \sqrt{\frac{\epsilon^2}{4} - \frac{1}{8} + \frac{1}{64\epsilon^2}} - \frac{1}{2}\right| \le \left|\frac{\epsilon}{2} + \frac{1}{2} + \frac{1}{8\epsilon} + \frac{\epsilon}{2} - \frac{1}{8\epsilon} - \frac{1}{2}\right| = \epsilon.$$
\newline

\textbf{2a.} Using limit laws, we have $\lim_{n \to \infty}(1 + \frac{n^2}{n^2 + 3})(3 + \frac{n}{n^2 + 6}) = \lim_{n \to \infty}(1 + \frac{n^2}{n^2 + 3}) \cdot \lim_{n \to \infty}(3 + \frac{n}{n^2 + 6}) = (1 + \lim_{n \to \infty}\frac{n^2}{n^2 + 3})(3 + \lim_{n \to \infty}\frac{n}{n^2 + 6})$.
\newline

Now, let's show that $\lim_{n\to\infty}\frac{n^2}{n^2 + 3} = 1$. We know that $|\frac{n^2}{n^2 + 3} - 1| = \frac{3}{n^2 + 3} < \epsilon$ if $\epsilon n^2 + 3n - 3 > 0$. Since it is an upward opening parabola in $n$, there exists $N$ such that $|\frac{n^2}{n^2 + 3} - 1| < \epsilon$.
\newline

Now, we show that $\lim_{n\to\infty}\frac{n}{n^2 + 6} = 0$. We can see that $\epsilon n^2 - n + 6\epsilon > 0$ implies $|\frac{n}{n^2 + 6}| < \epsilon$, and that is true for all $n$ greater than some $N$ because it represents an upward opening parabola.
\newline

Therefore, $(1 + \lim_{n \to \infty}\frac{n^2}{n^2 + 3})(3 + \lim_{n \to \infty}\frac{n}{n^2 + 6}) = (1 + 1)(3 + 0) = 6$.
\newline

\textbf{2b.} We see that $\lim_{n \to \infty} \frac{\sin^3 n + \cos n}{n^3 + 3n^2 + 1} = \lim_{n \to \infty} \frac{\sin^3 n}{n^3 + 3n^2 + 1} + \lim_{n \to \infty} \frac{\cos n}{n^3 + 3n^2 + 1}$. Also, $-\frac{1}{n^3 + 3n^2 + 1} \le \frac{\sin^3 n}{n^3 + 3n^2 + 1} \le \frac{1}{n^3 + 3n^2 + 1}$.
\newline

Now, let's show that $\frac{1}{n^3 + 3n^2 + 1}$ converges. We know that $|\frac{\sin^3 n}{n^3 + 3n^2 + 1} - 0| \le |\frac{1}{n^3 + 3n^2 + 1} - 0| < \epsilon$ if $\epsilon n^3 + 3\epsilon n^2 + \epsilon - 1 > 0$, which is true for all $n$ greater than some $N$ because it is a cubic with a positive leading coefficient.
\newline

Similarly for $\frac{\cos n}{n^3 + 3n^2 + 1}$, we get that $\lim_{n \to \infty} \frac{\sin^3 n + \cos n}{n^3 + 3n^2 + 1} = 0 + 0 = 0$.
\newline

\textbf{3.}
\newline

\textbf{4.} Assume for contradiction that $C < s$. This implies that $a_n < C < s$ for all $n$. If we set $\epsilon = s - C$, then $|a_n - s| = s - a_n < s - C$, which means $a_n > C$, which is a contradiction.
\newline

We can't conclude that $s < C$ because all the terms might tend to $C$ but never reach it. For example, consider the sequence $1 - \frac{1}{2^n}$. It converges to $1$, but the least upper bound is also $1$. So $s = C$.
\newline

For contradiction, assume that $C < s$. So $a_n \le C < s$ for all $n$. If $\epsilon = s - c$, then $|a_n - s| = s - a_n < s - C$, so $a_n > C$, which is a contradiction.
\newline

\textbf{5.} Call the sequence $a_n$. Since it converges (say to $s$), for all $\epsilon > 0$ there exists $N$ such that if $n > N$ then $|a_n - s| < \epsilon$. The values of $a_n$ for $n \le N$ form a finite set, which has a minimum and a maximum, say $A$ and $B$ respectively. For $n > N$, a lower bound is $s - \epsilon$ and an upper bound is $s + \epsilon$. So a lower bound of $a_n$ is $\min(A, s - \epsilon)$ and an upper bound is $\max(B, s + \epsilon)$.
\newline

An example of a bounded sequence that is not convergent is $(-1)^n$. It is always between $-1$ and $1$ but does not converge.
\newline

\textbf{6.} Let's construct a monotone subsequence, $b_n$, of $a_n$ assuming it is bounded. Define $b_1 = a_1$. If there are infinitely many elements in $a_n$ that are $\ge b_n$ for all $n$, let $b_{i + 1}$ be the first element after $b_i$ that is $\ge b_i$. If there aren't infinitely many elements in $a_n$ that are $\ge b_n$ for some $n$, then at one point there are infinitely many elements in $a_n$ that are less than that element. In that case, let $b_{i + 1}$ be the first element after $b_i$ that is $\le b_i$. This results in either a monotonically increasing or decreasing sequence.
\newline

\textbf{7.} Assume that a monotonically increasing sequence $a_n$ is bounded, and call the supremum $L$. By the definition of a supremum, for every $\epsilon > 0$, there exists $N$ such that $L - \epsilon < a_N \le L$. Since $a_n$ is increasing, for all $n \ge N$, we have $a_N \le a_n \le L$, so $L - \epsilon < a_n < L + \epsilon$, or equivalently $|a_n - L| < \epsilon$. The proof is similar for monotonically decreasing sequences.
\newline

Now, assume that a monotonically increasing sequence $a_n$ is convergent, and call the limit $L$. That means there exists $N$ such that for all $n > N$ we have $|a_n - L| < \epsilon$. Since $a_n$ is increasing, we have that $a_k < L + \epsilon$ for all $k \in \mathbb{N}$. Also, the lower bound is $a_1$ since it's increasing. The proof is similar for monotonically decreasing sequences.

\end{document}